% !TeX root = ../抽象代数笔记.tex
% 由主文档通过 \input 载入；请编译主文档。
\setcounter{chapter}{13} % 保留原书章号。
\chapter{域扩张与代数元素}\label{ch:extensions}

\begin{description}
  \item[学习目标] 把添加元素转化为域与向量空间的构造，使用最小多项式计算扩张次数。
  \item[先修知识] 第\ref{ch:polynomials}章，以及线性空间的基与维数。
\end{description}

\section{域扩张、生成域与扩张次数}
定义域扩张（field extension）$K/F$ 表示 $F$ 是 $K$ 的子域，将 $K$ 看作 $F$ 上的向量空间，并记 $[K:F]=\dim_FK$。定义生成域 $F(S)$ 与单扩张 $F(\alpha)$，区分它们与只允许有限次加法、乘法得到的子环 $F[S]$、$F[\alpha]$。

\section{代数元素与最小多项式}
定义代数元素（algebraic element）、超越元素（transcendental element）和首一最小多项式 $m_{\alpha,F}(x)$。用求值映射的核证明最小多项式存在时唯一且不可约，建立
\begin{equation}\label{eq:simple-extension}
  F(\alpha)=F[\alpha]\cong F[x]/(m_{\alpha,F}),
  \qquad [F(\alpha):F]=\deg m_{\alpha,F}.
\end{equation}
说明式\eqref{eq:simple-extension}的前提是 $\alpha$ 在 $F$ 上代数；超越情形对应有理函数域。

\section{塔式公式与基的构造}
对域塔 $F\subseteq K\subseteq L$，在 $[K:F]$ 和 $[L:K]$ 有限时，通过两组基的乘积构造证明 $[L:F]=[L:K][K:F]$。计算二次扩张与连续添加根的扩张，解释最小多项式会随基域改变，不能直接把每个已知多项式的次数相乘当作扩张次数。

\section{有限扩张与代数扩张}
定义代数扩张，证明有限扩张必代数、有限个代数元素生成有限扩张，以及代数性的传递性。区分“每个元素分别满足某个多项式”与“整个扩张具有有限维数”；代数而非有限的例子在第\ref{ch:splitting}章的代数闭包讨论中补充。

\section{例题、练习与自检}
计算 $\mathbb Q(\sqrt2,\sqrt3)$ 的基与次数，求 $\sqrt2+\sqrt3$ 的最小多项式，并计算扩域中元素的逆。练习要求先给出一个消去元素的多项式，再单独证明其不可约性；不能把任何以该元素为根的多项式都称为最小多项式。

